% ============================================================
% SECTION 2: BUSINESS REQUIREMENTS
% ============================================================
\section{Yêu cầu nghiệp vụ}
\label{sec:business-requirements}

Xuất phát từ nhu cầu thực tiễn của hệ thống theo dõi giao hàng và kỳ vọng của các bên liên quan (khách hàng, tài xế, quản trị viên), dự án xác định 8 yêu cầu nghiệp vụ cốt lõi. Nhằm tối ưu hóa kiến trúc hạ tầng, các yêu cầu này được phân loại thành hai nhóm chiến lược: Real-Time Processing (nhóm yêu cầu xử lý luồng thời gian thực) và Data Management \& Analytics (nhóm yêu cầu quản lý và phân tích dữ liệu).

% --- 2.1 Real-Time Processing Requirements ---
\subsection{Nhóm yêu cầu xử lý thời gian thực (Kafka Streams)}
\label{sec:realtime-requirements}

Trong bài toán theo dõi giao hàng thời gian thực, dữ liệu GPS từ tài xế phát sinh liên tục với tần suất cao (1 điểm/giây $\times$ hàng nghìn tài xế $\approx$ 86,4 triệu GPS points/ngày). Mô hình xử lý lô (batch processing) — chẳng hạn Hadoop MapReduce hay Spark Batch — gom dữ liệu theo chu kỳ (phút hoặc giờ) rồi mới xử lý, do đó \textbf{hoàn toàn không phù hợp} cho các tính năng đòi hỏi phản hồi tức thời: khách hàng cần thấy xe di chuyển mượt mà trên bản đồ, ETA cần cập nhật liên tục mỗi 10 giây, và cảnh báo tốc độ phải xuất hiện ngay khi vi phạm xảy ra. Thay vì gom lô, hệ thống cần chuyển sang mô hình \textit{``Dữ liệu chuyển động''} (Data in Motion) — xử lý ngay lập tức từng sự kiện GPS tại thời điểm nó phát sinh.

Apache Kafka Streams được chọn làm công nghệ xử lý cốt lõi vì ba lý do:
\begin{enumerate}
    \item Xử lý \textbf{từng sự kiện riêng lẻ} (\textit{event-at-a-time}) với độ trễ \textit{sub-millisecond} — đáp ứng yêu cầu cập nhật vị trí \textit{real-time} mà \textit{micro-batching} (Spark Streaming) không thể đạt được.
    \item Tích hợp sẵn \textbf{State Store} (RocksDB) cho xử lý có trạng thái — cho phép ``nhớ'' vị trí GPS trước đó để tính tốc độ mà không cần cơ sở dữ liệu bên ngoài (Redis, Memcached).
    \item Hỗ trợ \textbf{Windowed Aggregation} và \textbf{Stream-Table Join} — hai kỹ thuật cốt lõi cho việc tính tốc độ trung bình (ETA) và phát hiện khoảng cách đến điểm giao (geofencing) trong thời gian thực.
\end{enumerate}

Bốn yêu cầu xử lý thời gian thực cụ thể:

\begin{itemize}

% --- RT1 ---
    \item \textbf{Theo dõi vị trí thời gian thực (Real-Time Location Tracking)}
    \begin{itemize}
        \item \textbf{Input:} Luồng tọa độ GPS thô từ tài xế gửi đến Kafka topic \texttt{raw-location-events} (1 điểm/giây).
        \item \textbf{Output:} Dữ liệu vị trí đã được kiểm chứng (filtered/validated) gửi qua WebSocket lên giao diện khách hàng.
        \item \textbf{Công nghệ bắt buộc:} Kafka Streams -- cần lọc tọa độ không hợp lệ (lat/lng nằm ngoài phạm vi, accuracy > 20m) và \textit{deduplicate} các sự kiện gửi liên tục trong < 500ms. Điều này đòi hỏi \textbf{\textit{event-by-event processing}} (xử lý từng sự kiện), không thể dùng \textit{batch processing}.
        \item \textbf{Cassandra:} Bảng \texttt{trip\_locations} lưu trữ GPS trace phục vụ phát lại lịch sử.
    \end{itemize}

    % --- RT2 ---
    \item \textbf{Tính toán ETA động (Dynamic ETA Calculation)}
    \begin{itemize}
        \item \textbf{Input:} Luồng vị trí đã lọc + tọa độ điểm đích từ bảng \texttt{orders}.
        \item \textbf{Output:} ETA (giây) được tính lại mỗi 10 giây: $\text{ETA}_{\text{seconds}} = \frac{\text{distance\_km}}{\text{avg\_speed\_kmh}} \times 3600$.
        \item \textbf{Công nghệ bắt buộc:} Kafka Streams với \textbf{\textit{windowed aggregation}} (\textit{tumbling window} 30 giây) để tính tốc độ trung bình từ các GPS points gần nhất. Không thể tính đúng nếu dùng \textit{batch} vì \textit{batch} không đủ \textit{fresh} để cập nhật ETA liên tục.
    \end{itemize}

    % --- RT3 ---
    \item \textbf{Giám sát tốc độ \& cảnh báo (Speed Monitoring \& Alerts)}
    \begin{itemize}
        \item \textbf{Input:} Luồng vị trí đã lọc.
        \item \textbf{Output:} Cảnh báo gửi đến Kafka topic \texttt{alerts} khi tốc độ tức thời vượt ngưỡng 60 km/h.
        \item \textbf{Công nghệ bắt buộc:} Kafka Streams với \textbf{\textit{stateful processing}} -- sử dụng \textit{State Store} (RocksDB) để lưu vị trí GPS trước đó của mỗi tài xế, tính $v = \Delta d / \Delta t$. Chỉ có \textit{stream processing} mới so sánh được 2 điểm liên tiếp trong thời gian thực.
    \end{itemize}

    % --- RT4 ---
    \item \textbf{Cảnh báo khoảng cách - Geofencing (Proximity Alerts)}
    \begin{itemize}
        \item \textbf{Input:} Luồng vị trí + bảng tra cứu tọa độ điểm đích (KTable).
        \item \textbf{Output:} Cảnh báo khi tài xế cách điểm giao < 500m.
        \item \textbf{Công nghệ bắt buộc:} Kafka Streams \textbf{\textit{Stream-Table Join}} -- join luồng location với \textit{KTable} chứa tọa độ đích. Mỗi khi có \textit{location event} mới, hệ thống tra cứu điểm đích tương ứng và tính khoảng cách Haversine. Đây là pattern không thể thực hiện với \textit{batch processing}.
    \end{itemize}

\end{itemize}

% --- 2.2 Data Management & Analytics Requirements ---
\subsection{Nhóm yêu cầu quản lý dữ liệu \& phân tích (Cassandra)}
\label{sec:analytics-requirements}

Trong bài toán theo dõi giao hàng, việc lưu trữ toàn bộ dữ liệu lịch sử GPS và thông tin chuyến đi tạo ra một thách thức khổng lồ về mặt \textit{Volume} (khối lượng) và \textit{Velocity} (tốc độ ghi). Các hệ thống cơ sở dữ liệu quan hệ truyền thống (RDBMS) như MySQL hay PostgreSQL thường được thiết kế tối ưu cho các thao tác đọc (read-heavy) và đảm bảo tính toàn vẹn ACID, do đó bộc lộ nhiều điểm yếu chí mạng khi đối mặt với luồng dữ liệu IoT/GPS: nghẽn cổ chai khi số lượng giao dịch ghi (write transactions) tăng vọt, chi phí phần cứng đắt đỏ khi phải mở rộng dọc (vertical scaling), và sự phức tạp cùng rủi ro điểm lỗi đơn lẻ (\textit{Single Point of Failure} - SPOF) khi thực hiện phân mảnh dữ liệu (sharding). Để đáp ứng việc lưu trữ hàng tỷ bản ghi mỗi năm phục vụ truy xuất lịch sử và phân tích, hệ thống cần một giải pháp cơ sở dữ liệu phân tán thiết kế chuyên biệt cho thao tác ghi.

Apache Cassandra được chọn làm cơ sở dữ liệu cốt lõi cho nhóm yêu cầu này vì ba lý do chính:
\begin{enumerate}
    \item \textbf{Khả năng ghi tốc độ cực cao (\textit{High Write Throughput}):} Khác với RDBMS sử dụng B-Tree, Cassandra sử dụng cấu trúc \textit{Log-Structured Merge-Tree} (LSM-Tree) kết hợp \textit{Commit Log} và \textit{Memtable}, cho phép ghi dữ liệu tuần tự xuống đĩa với tốc độ cực nhanh mà không bị khóa (lock-free), cực kỳ phù hợp cho luồng dữ liệu \textit{write-heavy} như hệ thống GPS tracking.
    \item \textbf{Mở rộng ngang tuyến tính \& Không điểm lỗi đơn lẻ (\textit{Linear Scalability \& Masterless Architecture}):} Kiến trúc mạng lưới ngang hàng (\textit{peer-to-peer}) không có node chủ (master) giúp hệ thống phân tán đều dữ liệu và loại bỏ hoàn toàn SPOF. Việc tăng tải chỉ đơn giản là thêm node mới vào cụm (\textit{cluster}) mà không gây gián đoạn dịch vụ.
    \item \textbf{Mô hình dữ liệu tối ưu cho Time-Series (\textit{Time-Series Data Modeling}):} Cơ chế phân vùng (\textit{Partition Key}) và gom cụm (\textit{Clustering Key}) của Cassandra cho phép lưu trữ tập trung dữ liệu GPS của một chuyến đi trên cùng một node và tự động sắp xếp theo thứ tự thời gian, giúp các truy vấn lịch sử hành trình đạt hiệu suất đọc tối đa.
\end{enumerate}

Bốn yêu cầu quản lý dữ liệu và phân tích cụ thể:

\begin{itemize}
% --- DA5 ---
\item \textbf{Phát lại lịch sử hành trình (Trip History \& Playback)}
    \begin{itemize}
        \item \textbf{Bảng truy vấn:} \texttt{trip\_locations} — \textit{Partition Key} là \texttt{trip\_id} (để gom toàn bộ điểm GPS của một chuyến đi vào cùng một node), \textit{Clustering Key} là \texttt{timestamp DESC} (sắp xếp giảm dần theo thời gian).
        \item \textbf{Query:} \texttt{SELECT * FROM trip\_locations WHERE trip\_id = ?}
        \item \textbf{Lý do cần lưu trữ:} Dữ liệu quỹ đạo GPS (GPS trace) cần được giữ lại nguyên vẹn và vô thời hạn để phục vụ điều tra khiếu nại của khách hàng (dispute resolution), tra soát lộ trình thực tế và tuân thủ các quy định pháp lý về lưu trữ dữ liệu hành trình. Số lượng bản ghi cho mỗi chuyến đi có thể lên tới hàng nghìn điểm, đòi hỏi tốc độ ghi cực cao.
    \end{itemize}
% --- DA6 ---
\item \textbf{Phân tích hiệu suất tài xế (Driver Performance Analytics)}
   \begin{itemize}
        \item \textbf{Bảng truy vấn:} \texttt{driver\_analytics} — \textit{Partition Key} là \texttt{driver\_id}, \textit{Clustering Key} là \texttt{week\_start\_date DESC}. (Lưu ý: Các chỉ số được hệ thống Streaming tổng hợp sẵn và ghi xuống, vì Cassandra không tối ưu cho lệnh SUM/AVG trực tiếp).
        \item \textbf{Query:} Tổng hợp từ \texttt{trip\_metadata}: \texttt{SUM(total\_distance), AVG(average\_speed), COUNT(speeding\_violations)}.
        \item \textbf{Lý do cần lưu trữ:} Quản lý đội xe (Fleet Manager) cần theo dõi các chỉ số hiệu suất (KPIs) của tài xế theo từng tuần/tháng. Dữ liệu này là cơ sở để hệ thống đưa ra các chính sách thưởng/phạt tự động, xác định tài xế có rủi ro vi phạm an toàn cao và tối ưu hóa chất lượng dịch vụ.
    \end{itemize}   
% --- DA7 ---
\item \textbf{Tính cước phí giao hàng tự động (Automated Delivery Cost Calculation)}
    \begin{itemize}
        \item \textbf{Bảng truy vấn:} \texttt{trip\_metadata} — \textit{Partition Key} là \texttt{trip\_id}. Bảng này lưu trữ kết quả cuối cùng bao gồm tổng quãng đường (\texttt{distance\_km}), tổng thời gian (\texttt{duration\_min}). 
        \item \textbf{Công thức:} \texttt{cost = 3.00 + (0.50 × distance\_km) + (0.10 × duration\_min)}.
        \item \textbf{Truy vấn hiển thị:} \texttt{SELECT distance\_km, duration\_min, cost FROM trip\_metadata WHERE trip\_id = ?}
        \item \textbf{Lý do cần lưu trữ:} Cước phí chuyến đi phải minh bạch và được tính toán chuẩn xác dựa trên dữ liệu di chuyển thực tế thay vì mức giá ước lượng ban đầu. Việc lưu trữ \texttt{trip\_metadata} đảm bảo tính bất biến (immutable) của hóa đơn sau khi chuyến đi kết thúc, giúp cả tài xế và khách hàng có thể đối soát lại bất kỳ lúc nào.
    \end{itemize}
% --- DA8 ---
\item \textbf{Bản đồ nhiệt khu vực dịch vụ (Service Area Heatmaps)}
\begin{itemize}
    \item \textbf{Bảng truy vấn:} \texttt{trip\_locations} — query theo khoảng thời gian \texttt{timestamp}, nhóm theo geohash hoặc cluster tọa độ.
    \item \textbf{Query pattern:} \texttt{SELECT latitude, longitude FROM trip\_locations WHERE ...}.
    \item \textbf{Lý do cần lưu trữ:} Quản trị viên cần biết khu vực nào có nhu cầu cao để tối ưu phân bổ tài xế, và khu vực nào có ít đơn hàng để điều chỉnh chiến lược marketing.
\end{itemize}
    
\end{itemize}

\bigskip
\begin{table}[H]
\centering
\caption{Tổng hợp 8 yêu cầu nghiệp vụ}
\label{tab:requirements-summary}
\begin{tabular}{|p{1.5cm}|p{4cm}|p{3.5cm}|p{3.5cm}|}
\hline
\textbf{ID} & \textbf{Tên yêu cầu} & \textbf{Công nghệ chính} & \textbf{Output} \\
\hline
RT1 & Theo dõi vị trí real-time & Kafka Streams + Cassandra & WebSocket → Customer UI \\
\hline
RT2 & Tính ETA động & Kafka Streams (windowed agg) & ETA seconds → UI \\
\hline
RT3 & Giám sát tốc độ & Kafka Streams (stateful) & Alert → Admin Dashboard \\
\hline
RT4 & Geofencing alerts & Kafka Streams (KTable join) & Alert → Customer notification \\
\hline
DA5 & Phát lại hành trình & Cassandra & Route animation → Admin UI \\
\hline
DA6 & Phân tích hiệu suất & Cassandra & Dashboard charts → Admin UI \\
\hline
DA7 & Tính cước tự động & Cassandra & Trip cost → UI receipt \\
\hline
DA8 & Bản đồ nhiệt & Cassandra & Heatmap overlay → Admin UI \\
\hline
\end{tabular}
\end{table}
